% !TeX root = ../../Book3.tex
% 纲要标题：### 20.2 Serre 正规性判据
% 纲要位置：计划纲要.md，第 2592 行。

\section{Serre 正规性判据}
\label{b3:sec:2002}

分式在各高度一局部环中属于原环，还需要深度条件才能返回全环。Serre 判据把这一返回过程与一维局部环的 DVR 结构结合，给出适用于具有多个不可约分量的 Noether 环的正规性判别。


% 纲要内容（仅作写作计划，尚非教材正文）：
%
% * Noether 环上 R_1 与 S_2 的组合
% * 约化、全分式环整闭、局部环为整环的三步证明；以坐标幂等元排除不可约分量的局部相交
% * 正规局部整环的深度二与行列式技巧；汇总 Serre 条件、约化性与正规性的对应
%

\subsection{Noether 环上的 \texorpdfstring{\(R_1\)}{R1} 与 \texorpdfstring{\(S_2\)}{S2} 条件}
\label{b3:subsec:2002:01}

\begin{definition}\label{b3:def:normal-noether-ring}
设 \(R\) 为交换 Noether 环。若对每个素理想 \(\mathfrak p\)，\(R_{\mathfrak p}\) 都是在自身分式域中整闭的整环，则称 \(R\) 为\term[zhengguihuan]{正规环}[normal ring]。当 \(R\) 为整环时，这与 \(R\) 在 \(\operatorname{Frac}(R)\) 中整闭等价。
\end{definition}

本书的“正规环”采用这项逐点约定，允许原环有多个不可约分量；整环情形与分式域整闭性的等价已由\cref{b3:prop:normality-local-detection}证明。一般正规 Noether 环的不同不可约分量必须互不相交，而且各自是开闭子集，见下面的\cref{b3:cor:normal-product-domains}。例如，设 \(k\) 为域，则 \(k\times k\) 正规，而曲线 \(xy=0\) 在原点的两条局部分支相交，局部环 \(k[x,y]_{(x,y)}/(xy)\) 不正规。本节同时处理这种非整情形。

\ProseParagraphBreak
\(S_2\) 在判据中要解决的是清除分母。若 \(a\) 是非零因子，而 \(b\notin aR\)，非零商类 \(\bar b\in R/aR\) 必在某个关联素理想处仍然非零。下面两个引理说明，\(S_2\) 迫使这样的素理想高度为一。因此，一个分母无法清除的障碍不会只在高度至少二的素点处出现。

\begin{lemma}\label{b3:lem:s2-principal-associated-height}
设 \(R\) 为满足 \(S_2\) 的交换 Noether 环，\(a\in R\) 为非零因子且非单位。则
\[
\mathfrak p\in\operatorname{Ass}_R(R/aR)
\quad\Longrightarrow\quad\operatorname{ht}\mathfrak p=1.
\]
\end{lemma}
\begin{proof}
由\cref{b3:prop:associated-localization}，\(\mathfrak pR_{\mathfrak p}\) 是 \((R/aR)_{\mathfrak p}\) 的关联素理想，\cref{b3:cor:depth-zero-socle}给出这个非零局部模的深度为零。因为 \(a\in\mathfrak p\)，它在 \(R_{\mathfrak p}\) 上仍是非零因子且非单位。\cref{b3:cor:regular-quotient-depth}说明商去这个正则元素使深度降一，因而
\[
\operatorname{depth}R_{\mathfrak p}
=\operatorname{depth}_{R_{\mathfrak p}}(R/aR)_{\mathfrak p}+1=1.
\]
条件 \(S_2\) 要求 \(1\ge\min\{2,\dim R_{\mathfrak p}\}\)，故 \(\dim R_{\mathfrak p}\le1\)；深度不超过维数又使其至少为一。于是 \(\operatorname{ht}\mathfrak p=\dim R_{\mathfrak p}=1\)。
\end{proof}

\begin{lemma}\label{b3:lem:s2-fraction-membership}
设 \(R\) 为满足 \(S_2\) 的交换 Noether 环，\(a\in R\) 为非零因子，\(b\in R\)。若对每个高度一素理想 \(\mathfrak p\) 都有 \(b\in aR_{\mathfrak p}\)，则 \(b\in aR\)。
\end{lemma}
\begin{proof}
若 \(a\) 为单位，结论直接成立。否则反设 \(b\notin aR\)，则商类 \(\bar b\) 生成 \(R/aR\) 的一个非零有限生成子模，存在
\(\mathfrak p\in\operatorname{Ass}(R\bar b)\subseteq\operatorname{Ass}(R/aR)\)。上引理使 \(\operatorname{ht}\mathfrak p=1\)。在该处 \((R\bar b)_{\mathfrak p}\ne0\)，而它由 \(\bar b/1\) 生成，所以 \(\bar b/1\ne0\)。这正是 \(b/1\notin aR_{\mathfrak p}\)，与假设相矛盾。
\end{proof}

这一步把“某个分式不是环元素”变成“在某个高度一素点已经不是环元素”。关键是利用非零商类所生成子模的关联素理想；\(S_2\) 保证这个关联素理想恰好具有所需高度。\(R_1\) 随后提供这些高度一局部环的 DVR 结构，两项条件由此分别承担检测分母与清除分母的任务。

\subsection{局部环整闭性与 Serre 判据}
\label{b3:subsec:2002:02}

要从正规性推出 \(S_2\)，一维局部整环中的第一个非零因子很容易取得，困难在于局部维数至少二时再选第二个。商去第一个元素后，若极大理想成了关联素理想，就会得到一个不属于原环的分式。证明将按这个分式是否把极大理想映入自身分成两种情形：前者违背整闭性，后者会使极大理想成为主理想，违背维数假设。

\begin{lemma}\label{b3:lem:normal-local-depth-two}
设 \((R,\mathfrak m)\) 为交换 Noether 整闭局部整环。若 \(\dim R\ge2\)，则 \(\operatorname{depth}R\ge2\)。
\end{lemma}
\begin{proof}
选 \(0\ne a\in\mathfrak m\)，它是非零因子。只需证明
\(\mathfrak m\notin\operatorname{Ass}(R/aR)\)，然后用\cref{b3:lem:finite-prime-avoidance}在 \(\mathfrak m\) 中选第二个非零因子。

若反而存在 \(b\notin aR\) 使 \(\mathfrak mb\subseteq aR\)，置 \(z=b/a\in\operatorname{Frac}(R)\setminus R\)，则 \(z\mathfrak m\subseteq R\)。由于 \(\dim R\ge2\)，极大理想 \(\mathfrak m\ne0\)；又因 \(R\) 是整环，\(\operatorname{Ann}_R(\mathfrak m)=0\)，所以 \(\mathfrak m\) 是有限生成忠实模。假如 \(z\mathfrak m\subseteq\mathfrak m\)，乘以 \(z\) 就给出 \(\mathfrak m\) 的一个 \(R\) 线性自同态。这正使行列式技巧适用：一个分式稳定作用于这里的有限生成忠实理想时，其乘法自同态被首一多项式零化。对\cref{b3:thm:determinant-trick}取其中的理想 \(I=R\)，得到首一多项式 \(P(T)\in R[T]\)，使
\[
P(z)\mathfrak m=0.
\]
选取 \(0\ne v\in\mathfrak m\)，在分式域中由 \(P(z)v=0\) 得 \(P(z)=0\)。因此 \(z\) 对 \(R\) 整，整闭性使 \(z\in R\)，矛盾。

若 \(z\mathfrak m\) 不包含于 \(\mathfrak m\)，由于仍有 \(z\mathfrak m\subseteq R\)，便有 \(u\in\mathfrak m\) 使 \(zu\in R\setminus\mathfrak m\)，即 \(zu\) 为单位。对任意 \(v\in\mathfrak m\)，
\[
\frac vu=\frac{zv}{zu}\in R,
\]
故 \(\mathfrak m=uR\)。\cref{b3:thm:principal-ideal-theorem}使 \(\dim R\le1\)，又矛盾。因此 \(R/aR\) 的极大理想不是关联素理想，存在 \(c\in\mathfrak m\) 在 \(R/aR\) 上正则，\(a,c\) 即为长度二的正则序列。
\end{proof}

证明反向蕴含时，需要分别得到约化性、全分式环中的整闭性和每个局部环的整环性。\(R_0\) 与 \(S_1\) 先排除幂零元，高度一处的整闭性和 \(S_2\) 再把整分式追回原环。此时还没有得到局部环是整环；最后须用全分式环的坐标幂等元，把不同分量的分离带回原环，再利用局部性排除多个分量同时存在。

\begin{theorem}[Serre 正规性判据]\label{b3:thm:serre-normality}
设 \(R\) 为交换 Noether 环，不要求 \(R\) 为整环。则
\[
R\;\text{正规}\quad\Longleftrightarrow\quad R\;\text{同时满足}\;R_1\;\text{与}\;S_2.
\]
\end{theorem}
\begin{proof}
零环时两侧都是空条件；以下设 \(R\ne0\)。正规性由各素点的局部环定义，\(R_1\) 与 \(S_2\) 又由\cref{b3:prop:serre-conditions-localize}在局部化后保持，故可在任意给定素点处再局部化。

先设 \(R\) 正规。在高度零处局部环为域；在高度一处，一维 Noether 整闭局部整环由\cref{b3:thm:dvr-characterizations}为 DVR，所以 \(R_1\) 成立。对任意素点，若局部维数为一，整环中可取非零非单位，深度至少一；若局部维数至少二，上引理给出深度至少二；维数零处无要求。这证明 \(S_2\)。

反之，设 \(R_1\) 与 \(S_2\) 成立。它们分别蕴含 \(R_0\) 与 \(S_1\)，因此\cref{b3:prop:reduced-serre-criterion}给出 \(R\) 约化。

接着检验全分式环中的整闭性。取 \(Q(R)\) 中一个对 \(R\) 整的元素 \(z=b/a\)，其中 \(a\) 为非零因子。在每个高度一素理想 \(\mathfrak p\) 处，\(R_1\) 使 \(R_{\mathfrak p}\) 成为一维正则局部环。它由\cref{b3:prop:regular-local-domain}为整环，极大理想又由一个元素生成，故\cref{b3:thm:dvr-characterizations}使它成为 DVR，因而整闭。\(a/1\) 在这里仍为非零因子，分式 \(b/a\) 的像属于其分式域，并且仍满足同一个首一方程，所以属于 \(R_{\mathfrak p}\)。现在已在每个高度一处清除了分母，\cref{b3:lem:s2-fraction-membership}给出 \(b\in aR\)，即 \(z\in R\)。这证明 \(R\) 在 \(Q(R)\) 中整闭。

最后还须从全分式环的整闭性得到局部环是整环。把 \(R\) 局部化后，上述论证仍适用，因而现在可设 \(R\) 局部。\cref{b3:prop:reduced-total-quotient-product}给出
\[
Q(R)\cong\prod_{\mathfrak p\in\operatorname{Min}R}\operatorname{Frac}(R/\mathfrak p).
\]
若这个乘积有至少两个因子，则任一坐标幂等元 \(e_i=(0,\ldots,1,\ldots,0)\) 都不同于 \(0,1\)。它满足 \(e_i^2-e_i=0\)，对 \(R\) 整，刚证明的整闭性迫使 \(e_i\in R\)。但局部环只有幂等元 \(0,1\)：\(e_i+(1-e_i)=1\) 保证两者至少一个为单位，而 \(e_i(1-e_i)=0\) 随即迫使另一个为零，矛盾。因此上面的乘积只能有一个因子。约化性使全部极小素理想的交为零，所以唯一极小素理想就是零，\(R\) 是整环；已经证明的全分式环整闭性便成为分式域整闭性。各点局部环都如此，故原环正规。
\end{proof}

上述\term[Serre-zhenggui-xing-panju]{Serre 正规性判据}[Serre's criterion for normality]可与松村英之的证明比较\cite[Theorem 23.8, pp. 183--184]{Matsumura1989CommutativeRingTheory}。幂等元步骤把全分式环中的整闭性转化为定义所需的局部整环性，说明不同不可约分量在正规环中不能相交。

\subsection{非零因子与深度方法}
\label{b3:subsec:2002:03}

\begin{corollary}\label{b3:cor:normal-product-domains}
非零交换 Noether 环 \(R\) 正规，当且仅当它是有限个 Noether 整闭整环的直积。
\end{corollary}
\begin{proof}
若 \(R\) 正规，定理证明已说明它约化且在全分式环中整闭。全分式环的各坐标幂等元 \(e_1,\ldots,e_s\) 都属于 \(R\)，于是
\(R\cong Re_1\times\cdots\times Re_s\)。每个 \(Re_i\) 是 \(R\) 的商环，故为 Noether 环；它嵌入相应的分式域，所以是整环。若该分式域中的元素 \(\alpha\) 对 \(Re_i\) 整，把它放在第 \(i\) 个坐标、其余坐标置零，得到 \(\tilde\alpha\in Q(R)\)。将其首一方程的各低次系数看作 \(Re_i\subseteq R\) 的元素，最高次系数取 \(1_R\)，同一方程就在每个坐标中零化 \(\tilde\alpha\)。因此 \(\tilde\alpha\) 对 \(R\) 整，属于 \(R\)；又有 \(e_i\tilde\alpha=\tilde\alpha\)，故仍在 \(Re_i\) 中。反向由直积环的每个局部环都是一个因子的局部环，且整闭整环的局部化仍整闭得到。
\end{proof}

这里每个因子 \(Re_i\) 同时对应 \(\operatorname{Spec}R\) 中的 \(D(e_i)=V(1-e_i)\)，所以它既开又闭；因子为整环，又使这个子集不可约。这些开闭不可约分量互不相交并覆盖整个素谱。两张平面在一点相交的例子便违反了这一必要条件，即使它在所有高度一素点处都正则，仍不能正规。

\ProseParagraphBreak
各项局部条件与两个整体性质的关系汇总于\cref{b3:tab:serre-conditions}。

\begin{center}
\begin{minipage}{\textwidth}
\makeatletter
\def\@captype{table}
\makeatother
\centering
\caption[Serre 条件与约化性、正规性]{Serre 条件与约化性、正规性。设 \(R\) 为交换 Noether 环；每一行的左右两项等价，\(\mathfrak p\) 均遍历 \(R\) 的素理想。}
\label{b3:tab:serre-conditions}
\begin{tabularx}{\textwidth}{@{}>{\raggedright\arraybackslash}p{0.22\textwidth}>{\raggedright\arraybackslash}X@{}}
\toprule
条件 & 含义或等价性质\\
\midrule
\(R_0\) & 每个极小素理想处的局部环为域。\\
\(S_1\) & \(\operatorname{Ass}_R R=\operatorname{Min}R\)，即无嵌入关联素理想。\\
\(R_0\) 与 \(S_1\) & \(R\) 约化。\\
\(R_1\) & 高度不超过一处的局部环正则：高度零处为域，高度一处为 DVR。\\
\(S_2\) & 对每个 \(\mathfrak p\)，\(\operatorname{depth}R_{\mathfrak p}\ge\min\{2,\dim R_{\mathfrak p}\}\)。\\
\(R_1\) 与 \(S_2\) & \(R\) 正规。\\
\bottomrule
\end{tabularx}
\end{minipage}
\end{center}

表中 \(R_0\) 与 \(R_1\) 检验局部环的结构，\(S_1\) 与 \(S_2\) 检验深度；约化性和正规性都需要把这两类条件配合起来。

\ProseParagraphBreak
判据中的各个条件因而各有明确作用。Noether 性保证关联素理想有限、商模有可选的关联素理想，并使行列式技巧作用于有限生成极大理想。\(R_1\) 使高度一局部环成为整闭的 DVR。\(S_2\) 则保证分母清不掉的障碍已经出现在高度一处，而不是只藏在较高余维。

\ProseParagraphBreak
这也解释了为什么不能只凭“奇点集合余维至少二”就断言正规：该句话至多提供 \(R_1\)，还必须另有 \(S_2\)。例如 CM 环自动提供 \(S_2\)，于是对 CM 环，正规性确实等价于 \(R_1\)；\cref{b3:exm:r1-not-s2}中两张平面只在一点相交，正是缺少这项深度条件。正则环的所有局部环都正则，因而既满足 \(R_1\)，又由\cref{b3:cor:regular-ring-cm}为 CM 环、满足 \(S_2\)，所以正则环必正规。
